\section{\texorpdfstring{\Wjets}{W+jets} control region}%
\label{sec:WjetsCR}
The dominant background in the analysis arises from \Wjets background processes.
In order to effectively determine the background estimate and constrain
the normalisation of the \Wjets background across all measurement regions,
a dedicated control region is defined. This control region is designed to have
similar kinematic properties as the respective SRs but involves an inverted
variable selection. This inversion effectively reduces signal contamination
while increasing the yield and purity of a specific background.

The CR for the \Wjets background contains exactly three jets, which are selected
to ensure orthogonality to the SR.\@ At least one of the jets is \btagged.
Additionally the \Wboson-hadronic signal region is vetoed although the overlap is expected
to be small.
To enhance the fraction of \Wjets events, a cut centred around the \Wboson boson
mass is applied through the variable \mtW. Additionally, to reduce the \ttbar
contribution, a cut on \(\mblmin > \qty{180}{\GeV}\) is set.
The resulting \Wjets selection is shown in \cref{tab:RegionDefinitionsCR}
and the corresponding yields in \cref{tab:CR_Wsys:yields}.

\begin{table}[htbp]
  \begin{center}
    \begin{tabular}{l c}
      \toprule
      Selection              & \Wjets CR            \\
      \midrule
      \( N \)(central jet)   & \( = 3 \)            \\
      \( N_{\bjets} \)       & \( \geq 1 \)         \\
      \( \mtW + \met \)      & \( (60, \infty) \)   \\
      \( \pt(\ell) \)        & \( (30, \infty) \)   \\
      \( \pt(j)\)            & \( (25, \infty) \)   \\
      \( \mtW \)             & \( (40, 100) \)      \\
      \mblmin                & \( (180, \infty) \)  \\
      SR veto                & enabled              \\
      \bottomrule
    \end{tabular}
  \end{center}
  \caption{A summary of the event selection for the \Wjets control region.}%
  \label{tab:RegionDefinitionsCR}
\end{table}

\begin{table}[htbp]
  \begin{center}
    \input{tables/yields_CR_Wsys_l_3j_excl}
  \end{center}
  \caption{Yields of individual signal and background samples in the
    \Wjets control region.
    The total expected yield is using the nominal \tW DR sample.
    Uncertainties are statistical only.}%
  \label{tab:CR_Wsys:yields}
\end{table}


\cref{fig:WJetsCR:mblmin} shows the distribution of the \mblmin variable.
The \Wjets purity is roughly \qty{50}{\%} and there is no significant
sensitivity to the \ttbar and \tW interference observed.
A single-bin distribution is used in the likelihood fit as described
in \cref{sec:LikelihoodFit} and illustrated in \cref{fig:WJetsCR:nleptons}.

\begin{figure}[htbp]
  \centering
  \subfloat[]{
    \includegraphics[page=1, width=0.37\textwidth]{CR_WJets/CR_Wsys_l_3j_excl}%
    \label{fig:WJetsCR:mblmin}
  }
  \subfloat[]{
    \includegraphics[page=5, width=0.37\textwidth]{CR_WJets/CR_Wsys_l_3j_excl}%
    \label{fig:WJetsCR:nleptons}
  }

  \caption{Distributions in the \Wjets CR.\@
    \protect\subref{fig:WJetsCR:mblmin} \mblmin distribution and
    \protect\subref{fig:WJetsCR:nleptons} \(N_\mathrm{leptons} \) distribution.
    The full uncertainty without any fit taken into account is shown for the nominal setup using the \tW DR sample.
  }%
  \label{fig:WJetsCR}
\end{figure}
